\subsection{Centered-moment order}
\label{sec:moment-order}
\label{subsec:moment-monotonicity}

% ============================================================================
% BLURB A. Generator on centered monomials
%
% Reader question:
%   Why should the semigroup preserve the centered-moment order?
%
% Job:
%   Start directly from the generator acting on
%       \chi_A^*=\prod_{j\in A}(\eta(j)-p_j^\star).
%   For i\in A, first record the one-site identity
%       \lambda_i(E_{i,\rho_i}-I)\chi_A^*
%       = \lambda_i(\rho_i-p_i^\star)\chi_{A\setminus\{i\}}^*
%         - \lambda_i\chi_A^*.
%   Then expand \lambda_i(\rho_i-p_i^\star) and \lambda_i in ordinary
%   monomials. This gives the central generator formula
%       L\chi_A^*
%       = sum_{i\in A} [
%           -\lambda_i(0)\chi_A^*
%           + sum_{S\subseteq N(i)}
%             (\widehat{\lambda_i\rho_i}(S)
%              -p_i^\star\widehat{\lambda_i}(S))
%             \chi_S\chi_{A\setminus\{i\}}^*
%           + sum_{\emptyset\ne S\subseteq N(i)}
%             (-\widehat{\lambda_i}(S))\chi_S\chi_A^*
%         ].
%
% Reader takeaway:
%   Preservation of centered moments is a direct generator statement. No patch
%   representation is needed in this part of the proof.
%
% Keep out:
%   Do not introduce the old auxiliary notation h_i(S), b_i(S), or r_i.
% ============================================================================

% ============================================================================
% BLURB B. Why the coefficients are nonnegative
%
% Reader question:
%   Where does patch positivity enter the centered-basis generator calculation?
%
% Job:
%   Use part (i) exactly where it is needed. For nonempty S,
%       -\widehat{\lambda_i}(S) >= 0.
%   For the other coefficient family, show
%       \widehat{\lambda_i\rho_i}(S)
%       -p_i^\star\widehat{\lambda_i}(S) >= 0
%   for every S\subseteq N(i). For S=\emptyset this is
%       \lambda_i(0)(\rho_i(0)-p_i^\star) >= 0,
%   using p_i^\star <= \rho_i(0). For nonempty S it follows from the threshold
%   formula together with the coefficient criterion.
%
%   Then explain why multiplying by an ordinary monomial does not spoil
%   positivity in the centered basis. Use the sitewise identities
%       \eta(j)=p_j^\star+(\eta(j)-p_j^\star),
%       \eta(j)(\eta(j)-p_j^\star)
%       =(1-p_j^\star)\eta(j),
%   and 0<=p_j^\star<=1 to conclude that every
%       \chi_S\chi_B^*
%   is a nonnegative linear combination of centered monomials.
%
% Reader takeaway:
%   The threshold p^\star is exactly the centering for which all off-diagonal
%   coefficients of the generator become nonnegative.
% ============================================================================

% ============================================================================
% BLURB C. From the generator to preservation of order
%
% Reader question:
%   How do the coefficient signs imply the semigroup comparison?
%
% Job:
%   Work first in a finite zero-boundary volume. The generator matrix in the
%   centered-monomial basis has nonnegative off-diagonal entries, so its
%   exponential preserves the nonnegative orthant of centered-moment vectors.
%   Apply this to the difference of two initial laws to obtain
%       \mu\preceq_*\nu
%       => \mu P_t\preceq_*\nu P_t
%   in finite volume.
%
%   Then pass to Lambda using the finite-propagation approximation from the
%   appendix. Keep this passage short and explicit; do not re-prove finite
%   propagation here.
%
% Reader takeaway:
%   This proves the main assertion of Theorem B(ii).
% ============================================================================

% ============================================================================
% BLURB D. Immediate consequences and completion of Theorem B
%
% Reader question:
%   How do the remaining assertions of part (ii) follow?
%
% Job:
%   Apply the same cone preservation to a single initial law to conclude
%       \mu\in\mathcal M_* => \mu P_t\in\mathcal M_*.
%   Then derive ordinary monomial comparison from
%       \chi_A
%       = sum_{B\subseteq A}
%         (prod_{i\in A\setminus B}p_i^\star)\chi_B^*,
%   whose coefficients are nonnegative. Do not use the patch representation
%   for this consequence.
%
%   Finally note that every nonempty centered moment of
%   \mu_{\mb p^\star} is zero, hence
%       \mu_{\mb p^\star}\preceq_*\mu
%   for every \mu\in\mathcal M_*, giving the final comparison in the theorem.
%   End explicitly by saying that part (ii), and therefore
%   Theorem~\ref{thm:patch-positivity-order}, is proved.
%
% Keep out:
%   Do not restore the old product-profile corollary here unless it is later
%   judged independently useful.
% ============================================================================

\subsection{Pure-death comparison}
\label{subsec:pure-death-comparison}

% Theorem~\ref{thm:patch-positivity-order} is complete before this subsection
% begins. This is the first place where the centered-moment order and the patch
% representation need to be combined.
%
% Main job:
%   Introduce the centered expansion of end factors here, because it is needed
%   for the comparison of two patch representations. For every end patch P
%   based at i, use affinity and the definition of p_i^star to write
%       C(z,P)
%       = C(p_i^star,P) + kappa(P)(z-p_i^star),
%   with C(p_i^star,P) >= 0 and kappa(P) >= 0. Expanding the product over end
%   patches expresses its expectation as a sum of centered moments with
%   nonnegative coefficients. Hence it is nonnegative for mu in M_*.
%
%   Then prove the pure-death comparison itself: removing the
%   environment-independent 1->0 transitions preserves patch positivity and the
%   threshold profile, and increases the relevant bulk and end patch
%   contributions. The centered end-factor expansion allows the comparison to
%   be made term by term for every mu in M_* and every successful-interaction
%   skeleton before averaging.
%
%   Put the invariant-measure comparison here afterwards as a consequence.
%
% Keep out:
%   Do not present the centered end-factor expansion as part of the proof of
%   centered-moment order; that theorem has already been proved by the generator
%   calculation.
%
% Placeholder. Rewrite subsection blurb by blurb.
